% Advanced poetry branch. P means project-composed; R means received text.
% Authoritative received-text windows reused inside the advanced poetry branch.
% Callers supply the source label, apparatus, and surrounding quote environment.

\newcommand{\PoetryPhaedrusPrologueOpening}{%
  \Latin{Aesopus auctor quam materiam repperit,\\
  hanc ego polivi versibus senariis.}}

\newcommand{\PoetryVirgilAeneidOpening}{%
  \Latin{Arma virumque cano, Troiae qui primus ab oris\\
  Italiam, fato profugus, Lavinaque venit\\
  litora, multum ille et terris iactatus et alto\\
  vi superum, saevae memorem Iunonis ob iram.}}

\newcommand{\PoetryHoraceCarminaOneEleven}{%
  \Latin{Tu ne quaesieris---scire nefas---quem mihi, quem tibi\\
  finem di dederint, Leuconoe, nec Babylonios\\
  temptaris numeros. Ut melius, quidquid erit, pati,\\
  seu plures hiemes seu tribuit Iuppiter ultimam,\\
  quae nunc oppositis debilitat pumicibus mare\\
  Tyrrhenum: sapias, vina liques, et spatio brevi\\
  spem longam reseces. Dum loquimur, fugerit invida\\
  aetas: carpe diem, quam minimum credula postero.}}

% The Missal prints the received wording as compact paragraphs.  P6 and P7
% display each three-member stanza as three analytical lines; the lineation is
% course apparatus and is declared beside each use.
\newcommand{\PoetryPentecostSequenceComplete}{%
  \Latin{Veni, Sancte Spiritus,\\
  et emitte cælitus\\
  lucis tuæ radium.\\[0.2\baselineskip]
  Veni, pater pauperum,\\
  veni, dator munerum,\\
  veni, lumen cordium.\\[0.2\baselineskip]
  Consolator optime,\\
  dulcis hospes animæ,\\
  dulce refrigerium.\\[0.2\baselineskip]
  In labore requies,\\
  in æstu temperies,\\
  in fletu solacium.\\[0.2\baselineskip]
  O lux beatissima,\\
  reple cordis intima\\
  tuorum fidelium.\\[0.2\baselineskip]
  Sine tuo numine,\\
  nihil est in homine,\\
  nihil est innoxium.\\[0.2\baselineskip]
  Lava quod est sordidum,\\
  riga quod est aridum,\\
  sana quod est saucium.\\[0.2\baselineskip]
  Flecte quod est rigidum,\\
  fove quod est frigidum,\\
  rege quod est devium.\\[0.2\baselineskip]
  Da tuis fidelibus,\\
  in te confidentibus,\\
  sacrum septenarium.\\[0.2\baselineskip]
  Da virtutis meritum,\\
  da salutis exitum,\\
  da perenne gaudium. Amen. Alleluia.}}


\begin{courseunit}{P1}
\chapter{Poetry branch: textual prosody}

\textbf{Starting point.} Begin this branch after M22, when poetic
order can be recovered without depending on line order and the full prose
grammar is reasonably dependable.

\textbf{R, PHAED-POST1919, Phaedrus 1.1.1--2:}
\begin{quote}
\Latin{Ad rivum eundem lupus et agnus venerant,\\
siti compulsi. Superior stabat lupus.}
\end{quote}

J.~P. Postgate, ed., \textit{Phaedri Fabulae Aesopiae} (Oxford: Clarendon,
1919), printed p.~1, scan PDF p.~35, visually checked; \Latin{u/v} is
normalized. Project translation: \English{A wolf and a lamb had come to the
same stream, driven by thirst. The wolf stood higher up.}

The same checked page supplies the exact metrical claim used below:

\textbf{R, PHAED-POST1919, Book 1, Prologue 1--2:}
\begin{quote}
\PoetryPhaedrusPrologueOpening
\end{quote}

The course normalizes \Latin{u/v}. Project translation: \English{Aesop, the
author, found the subject matter; I have polished it in senarian verses.}

\section{The senarius working frame}

Bennett, \textit{New Latin Grammar} \S370.1, identifies the Latin iambic
trimeter also as the \emph{senarius}: six iambic feet. Its pure frame, with
the course's ordinary final anceps notation, is

\[
\smile\,\text{--}\mid
\smile\,\text{--}\;\Vert\;
\smile\,\text{--}\mid
\smile\,\text{--}\;\Vert\;
\smile\,\text{--}\mid
\smile\,\times .
\]

The double bars group the six feet into three dipodies; they are not caesura
marks.
This frame is a starting contract, not permission to force every position
into a pure iamb. Bennett \S370.2 permits a tribrach in any foot except the
last and, in the odd-numbered feet, a spondee, dactyl, or anapaest. The caesura
normally falls within the third foot and less often within the fourth. Write
the six-foot frame first; then record every proposed substitution and its
quantity evidence. The Phaedrus prologue tells us which family to test, while
the fable line decides whether a particular test survives.

\section{Quantity is textual evidence}

A syllable may be long by nature, long by position, or short. Record nature
from a lexicon or the supplied memory sheet; do not infer it from English.
\Latin{ae}, \Latin{au}, and \Latin{oe} normally form long diphthongs. A vowel
before two consonants is normally long by position for meter, although a stop
plus liquid deserves separate testing. A vowel or vowel plus final
\Latin{-m} before a following vowel or \Latin{h-} is normally available for
elision. The last syllable of a line is metrically indifferent.

Mark syllable boundaries first, then natural quantities, then positional
length, then possible elision. Keep three marks distinct: known, inferred,
and unresolved. Only after that compare the line with a metrical scheme.
The printed prologue identifies the expected medium as \Latin{versibus
senariis}; that is evidence for which scheme to test, not a completed
scansion. Test the fable lines as iambic senarii only after recording their
syllables and quantities, and leave any unresolved position visible.

From memory write the senarius frame and its substitution positions, the
diphthongs, positional rule, elision environment, and meaning of anceps. Parse
\Latin{eundem}, \Latin{venerant}, \Latin{compulsi}, and \Latin{superior}; then
mark the text without translating. Keep a scansion ledger with columns for
syllable, quantity, reason, position, and uncertainty.

\ifsolutions
\section{Model and points to check}
The prose spine is \Latin{lupus et agnus venerant; lupus stabat}.
\Latin{siti compulsi} agrees with both animals and gives the cause of arrival;
\Latin{superior} prepares the fable's spatial and moral asymmetry. A quantity
mark without a stated reason is a prompt for checking, not yet an answer.
\else
\section{Before continuing}
You should reasonably be able to syllabify a fresh line, mark secure natural
and positional quantities, identify possible elision, test rather than assume
a scheme, and keep unresolved quantities visible.
\fi
\end{courseunit}

\begin{courseunit}{P2}
\chapter{Poetry branch: hexameter}

\textbf{R, VERG-HEND1891, \textit{Aeneid} 1.1--4:}
\begin{quote}
\PoetryVirgilAeneidOpening
\end{quote}

John Henderson, ed., \textit{Aeneid, Book I} (Toronto: Copp, Clark, 1891),
printed p.~27, scan PDF p.~31, visually checked. The edition prints
\Latin{Trojae} and \Latin{Lavinaque}; the course normalizes \Latin{j/i} but
retains its \Latin{Lavinaque} reading. Project translation: \English{Arms and
the man I sing, who first from the shores of Troy, an exile by fate, came to
Italy and the Lavinian shores, much buffeted on lands and on the deep by the
power of the gods above, because of the mindful anger of savage Juno.}

\section{Six feet and one sentence}

A hexameter has six feet. Feet one through four may be dactyls
(long--short--short) or spondees (long--long); the fifth is normally a dactyl;
the sixth is long plus anceps. Word boundary and foot boundary are different.
A principal caesura often gives the long line an internal turn, but syntax may
continue far beyond it.

The first line tests as
\[
\text{-- $\smile\smile$ | -- $\smile\smile$ | -- -- | -- -- |
-- $\smile\smile$ | -- $\times$},
\]
with word groups \Latin{arma vi | rumque ca | no Tro | iae qui | primus ab |
oris}. The display is a metrical model; ordinary spelling remains the source
text. Verify each length on the quantity sheet, including consonant groups.

Recover the delayed objects of \Latin{cano}, the antecedent of \Latin{qui},
the destinations governed by \Latin{venit}, and the forces expressed by
\Latin{fato}, \Latin{vi}, and \Latin{ob iram}. Then scan two further lines and
write a prose-order paraphrase labeled P.

\ifsolutions
\section{Model and points to check}
The opening postpones the full subject of \Latin{venit} while accumulating
origin, destination, fate, wandering, and divine cause. A prose-order check may
begin \Latin{cano arma et virum qui, profugus fato ab oris Troiae, primus
Italiam Lavinaque litora venit}. It is not a variant of Virgil.
\else
\section{Before continuing}
You should reasonably be able to scan and construe a fresh hexameter, keep foot
and word boundaries distinct, locate a plausible caesura, and explain how
metrical order serves a sentence extending through several lines.
\fi
\end{courseunit}

\begin{courseunit}{P3}
\chapter{Poetry branch: elegiac verse}

\textbf{R, OVID-OWEN1885, \textit{Tristia} 1.1.1--2:}
\begin{quote}
\Latin{Parve---nec invideo---sine me, liber, ibis in urbem:\\
ei mihi, quod domino non licet ire tuo!}
\end{quote}

S.~G. Owen, ed., \textit{P. Ovidi Nasonis Tristium libri V} (Oxford:
Clarendon, 1885), opening text leaf of book 1, scan PDF p.~70, visually checked.
Project translation: \English{Little book---I do not begrudge it---you will go
to the city without me: alas for me, because your master is not permitted to
go!}

\section{Couplet and turn}

An elegiac couplet joins a hexameter to a pentameter. Scan the hexameter by the
P2 method. The pentameter consists of two balanced halves around a fixed
division: two-and-a-half feet, then two-and-a-half feet. The second half ends
with two dactyls and a final long or anceps syllable. No word normally bridges
the central division.

Before scansion, mark the address \Latin{parve liber}, parenthetic
\Latin{nec invideo}, future \Latin{ibis}, impersonal \Latin{non licet}, and
dative \Latin{domino tuo}. Then describe the turn: the hexameter grants the
book a journey; the pentameter measures that freedom against the speaker's
exclusion.

Write two P couplet-shaped prose pairs, each with a clear turn after the first
line. Only then attempt a bounded metrical couplet, keeping a line-by-line
quantity ledger and changing the thought rather than corrupting quantity when
a word will not fit.

\ifsolutions
\section{Model and points to check}
The diminutive \Latin{parve}, direct address, and future create intimacy; the
impersonal construction shifts from what the book will do to what is denied
its master. A sound exercise may remain an accurate prose pair until its
quantities have been independently checked.
\else
\section{Before continuing}
You should reasonably be able to distinguish hexameter from pentameter, locate
the pentameter division, explain a couplet's semantic turn, and revise a
bounded attempt from an explicit quantity ledger.
\fi
\end{courseunit}

\begin{courseunit}{P4}
\chapter{Poetry branch: lyric measure and concentrated counsel}

\textbf{R, HOR-WICKGAR1912, Horace, \textit{Carmina} 1.11:}
\begin{quote}
\PoetryHoraceCarminaOneEleven
\end{quote}

E.~C. Wickham and H.~W. Garrod, eds., \textit{Q. Horati Flacci opera}
(Oxford: Clarendon, 1912), poem 1.11 on scan PDF p.~29, visually checked. The
course expands the edition's initial small-cap \Latin{TV}.

\textbf{Project translation.} Do not ask---it is forbidden to know---what end
the gods have given me, what end to you, Leuconoe; and do not test Babylonian
reckonings. How much better to endure whatever shall be, whether Jupiter has
granted more winters or this as the last, which now weakens the Tyrrhenian sea
against opposing rocks. Be wise, strain the wine, and cut back long hope within
a brief span. While we speak, envious time will have fled: seize the day,
trusting tomorrow as little as possible.

\section{A lyric line is a repeated contract}

This ode uses the greater Asclepiadean line. For a working choriambic scan, its
full repeated quantity scheme is:

\[
\underbrace{\text{--}\;\text{--}}_{\text{opening base}}
\mid
\underbrace{\text{--}\,\smile\smile\,\text{--}}_{\text{choriamb 1}}
\parallel
\underbrace{\text{--}\,\smile\smile\,\text{--}}_{\text{choriamb 2}}
\parallel
\underbrace{\text{--}\,\smile\smile\,\text{--}}_{\text{choriamb 3}}
\mid
\smile\,\times
\]

Thus the line has three, not two, choriambic sequences
(long--short--short--long), with stable internal divisions. Allen and
Greenough, \textit{New Latin Grammar} \S\S625.6 and 626.8 (printed pp.~420 and
422; scan PDF pp.~438 and 440, visually checked), give the exact line and
identify Horace's Ode 1.11 as the major Asclepiadean system. Their preceding
note explains that ``choriambic'' is a useful later way of scanning a
logaoedic form; the working display above is not a complete historical theory
of the meter. Copy the scheme above one line, then test every later line
against the same template; never invent a new template to rescue an error.

Construe before interpreting: prohibit \Latin{quaesieris} and
\Latin{temptaris}; attach both indirect questions to \Latin{scire}; take
\Latin{quidquid erit} with \Latin{pati}; distinguish the alternatives after
\Latin{seu}; then follow the jussive sequence \Latin{sapias, liques, reseces}.
The final imperatives concentrate the whole meditation.

Compare this long, internally articulated lyric line with the project-built
hymn-style stanza printed in the branch exercises. Record quantitative
scheme, clause length, repetition, and source attribution separately; a
similar devotional effect does not establish a shared meter or author.

\ifsolutions
\section{Model and points to check}
The movement is prohibition, alternative futures, practical counsel, and final
maxim. \Latin{fugerit} looks ahead from the ongoing \Latin{loquimur}; the poem
makes time vanish during the act of speech. A defensible lyric reading joins
that grammar to the repeated line form rather than treating
\Latin{carpe diem} as an isolated slogan.
\else
\section{Before continuing}
You should reasonably be able to test a lyric line against a fixed scheme,
construe Horace's subordinate architecture, compare lyric and hymn forms
without collapsing them, and explain how a stanza concentrates an argument.
Test that ability on a second line before continuing: write the unchanged
scheme first, divide and justify the quantities, recover every grammatical
dependency, and distinguish observation from interpretation. If the scheme is
being altered merely to fit the line, return to the complete Horace text and
the greater-Asclepiadean memory sheet. If the scan is secure but the sentence
remains obscure, make a finite-verb and clause ledger before trying another
translation. Continue when both the fixed form and the movement of thought can
be explained from the printed evidence.
\fi
\end{courseunit}

\begin{courseunit}{P5}
\chapter{Poetry branch: Christian poetics}

\textbf{R, MR62-FORM-PENT, Pentecost Sequence:}
\begin{quote}
\Latin{Consolator optime, dulcis hospes animæ, dulce refrigerium.\\
In labore requies, in æstu temperies, in fletu solacium.\\
O lux beatissima, reple cordis intima tuorum fidelium.\\
Sine tuo numine, nihil est in homine, nihil est innoxium.\\
Lava quod est sordidum, riga quod est aridum, sana quod est saucium.\\
Flecte quod est rigidum, fove quod est frigidum, rege quod est devium.}
\end{quote}

\textit{Missale Romanum} (1962), Pentecost, printed p.~357, local evidence PDF
p.~438, scan leaf 437, visually checked. No author claim is attached here; a
liturgical witness verifies wording and use, not a debated attribution.

\section{Read inheritance and transformation}

Christian poetry may employ classical quantity, late antique forms, scriptural
language, syllabic or rhythmic organization, rhyme, liturgical recurrence, or
several of these at once. Date, meter, genre, use, and authorship are separate
questions. Record each claim with its own evidence.

The sequence builds by paired nouns and imperatives. Mark the three situational
pairs in \Latin{in labore, in æstu, in fletu}; the privative pair
\Latin{sine tuo numine / nihil}; and the two triads
\Latin{lava, riga, sana} and \Latin{flecte, fove, rege}. Parse every relative
\Latin{quod}; its neuter antecedent is supplied by the predicative adjective,
not by a preceding expressed noun.

Make a lexical field table: need or wound, divine action, desired result,
scriptural resonance, and whether that resonance has been verified. Then
compare the poem's petitions with the prose Collect and Postcommunion in
\texttt{MR62-FORM-PENT}. Composition should imitate one structural feature at
a time and remain labeled P.

\ifsolutions
\section{Model and points to check}
The first two lines name the addressee through apposition and paradoxical
relief; the next two turn to the interior person; the final two triads move
from cleansing through restoration to guidance. A careful note may say
\English{the wording recalls biblical images}; an exact allusion claim waits
for a cited source-context window.
\else
\section{Before continuing}
You should reasonably be able to construe a Christian poem as poetry, separate
authorship from liturgical witness, trace lexical and structural development,
and compare verse petition with prose prayer from exact evidence.
\fi
\end{courseunit}

\begin{courseunit}{P6}
\chapter{Poetry branch: rhythmic hymnody and sequences}

\section{The two received comparison texts}

\textbf{R/N, MR62-FORM-PENT, complete Pentecost Sequence:} the wording is
received from the 1962 \textit{Missale Romanum}, printed p.~357, local evidence
PDF p.~438, scan leaf 437, visually checked. The course omits the printed
accents, regularizes \Latin{u/v}, and displays each three-member stanza as
three analytical lines. Those line breaks are course apparatus, not a claim
about performance or the Missal's compact typography.

\begin{quote}
\small
\PoetryPentecostSequenceComplete
\end{quote}

\textbf{R, VERG-HEND1891, \textit{Aeneid} 1.1--4:} John Henderson, ed.,
\textit{Aeneid, Book I} (Toronto: Copp, Clark, 1891), printed p.~27, scan PDF
p.~31, visually checked. The edition prints \Latin{Trojae} and
\Latin{Lavinaque}; the course normalizes \Latin{j/i} but retains its
\Latin{Lavinaque} reading. This is the complete four-line window used here.

\begin{quote}
\PoetryVirgilAeneidOpening
\end{quote}

Using the complete sequence above, copy its analytical lines in three-member
groups and make three independent records:

\begin{enumerate}
\item a grammar record of finite verbs, vocatives, ellipsis, and parallel
constructions;
\item a visible-pattern record of syllable counts, three-member grouping,
parallel syntax, repeated words, and written line endings;
\item a quantity record made by the P1 method, including uncertainty.
\end{enumerate}

\section{Do not turn one record into another}

\textbf{Quantitative structure} organizes long and short syllables by a scheme.
\textbf{Page-based textual pattern} records recurring syllable counts, line
groups, parallel constructions, repeated words, and written endings.
\textbf{Rhyme and assonance} are claims about sound; visible spelling may
suggest a question but cannot establish its spoken answer in a course that
supplies no pronunciation system. Keep the quantity, grammar, page pattern,
and written-ending records separate.

Test the visible pairs \Latin{radium/cordium},
\Latin{refrigerium/solacium}, \Latin{fidelium/innoxium},
\Latin{saucium/devium}, and \Latin{septenarium/gaudium}. Note exact and inexact
correspondence. Then compare the sequence's paired development with a hymn
project-built stanza printed later in this packet's integrated practice and
with the complete Virgilian sentence printed above. Explain what changes when
the line is a self-contained petition rather than one segment of an epic
period.

\section{Reconstruction and bounded imitation}

Reconstruct the sequence from its imperative chain:
\Latin{veni---emitte---veni---veni---veni---reple---lava---riga---sana---
flecte---fove---rege---da---da---da---da}. The four occurrences each of
\Latin{veni} and \Latin{da} are part of the poem's large architecture, not
dispensable duplicates. Then write four P lines with a repeated grammatical
pattern and two exact end correspondences. First secure grammar and meaning;
only a later draft may adjust rhythm. Record every sacrificed or substituted
word.

\ifsolutions
\section{Model and points to check}
The imperative ledger reveals large-scale form even before metrical judgment.
A sound comparison can say that Virgil postpones completion across lines while
the sequence often closes a petition or balanced member at a line boundary.
Neither observation depends on a reconstructed performance.
\else
\section{Before continuing}
You should reasonably be able to distinguish quantitative evidence from
grammar, syllable count, line grouping, and written-ending evidence; analyze a
complete sequence; reconstruct its verbal architecture; and create a bounded
patterned imitation with a revision record.
\fi
\end{courseunit}

\begin{courseunit}{P7}
\chapter{Poetry branch: source-first commentary}

\textbf{R, PHAED-POST1919, Phaedrus 1.1.1--6:}
\begin{quote}
\Latin{Ad rivum eundem lupus et agnus venerant,\\
siti compulsi. Superior stabat lupus,\\
longeque inferior agnus. Tunc fauce improba\\
latro incitatus iurgii causam intulit;\\
Cur, inquit, turbulentam fecisti mihi\\
aquam bibenti? Laniger contra timens \ldots}
\end{quote}

J.~P. Postgate, ed., \textit{Phaedri Fabulae Aesopiae} (Oxford: Clarendon,
1919), printed p.~1, scan PDF p.~35, visually checked; \Latin{u/v} is
normalized. This expands the P1 window on the same verified printed page. The
final ellipsis marks the deliberate boundary after line 6, not missing text in
the witness.

\section{The comparison dossier is printed here}

The three received texts used by this packet's comparative work follow.
They reuse the exact verified windows and normalization decisions identified
in the owner passage inventory.

\textbf{R, VERG-HEND1891, Virgil, \textit{Aeneid} 1.1--4:}
\begin{quote}
\PoetryVirgilAeneidOpening
\end{quote}
John Henderson, ed., \textit{Aeneid, Book I} (Toronto: Copp, Clark, 1891),
printed p.~27, scan PDF p.~31, visually checked. This is the complete
four-line window; the course normalizes the edition's \Latin{j/i} and retains
its \Latin{Lavinaque} reading.

\textbf{R, HOR-WICKGAR1912, Horace, \textit{Carmina} 1.11:}
\begin{quote}
\PoetryHoraceCarminaOneEleven
\end{quote}
E.~C. Wickham and H.~W. Garrod, eds., \textit{Q. Horati Flacci opera}
(Oxford: Clarendon, 1912), scan PDF p.~29, visually checked. The complete ode
is printed here; no material variant or further normalization is supplied.

\textbf{R/N, MR62-FORM-PENT, complete Pentecost Sequence:} 1962
\textit{Missale Romanum}, printed p.~357, local evidence PDF p.~438, scan leaf
437, visually checked. The complete sequence is printed here; accents are
omitted, \Latin{u/v} is regularized, and the three-member analytical lineation
is course apparatus rather than received typography.
\begin{quote}
\small
\PoetryPentecostSequenceComplete
\end{quote}

\section{Seven layers, kept distinct}

A compact commentary should proceed in this order:

\begin{enumerate}
\item witness, edition, passage boundary, and normalization;
\item text, declared normalization, and any supplied variant that materially
affects the reading (or an explicit note that none is supplied);
\item meter or visible pattern, with quantity reasons where applicable and
unresolved points kept visible;
\item syntax and a temporary prose order;
\item diction, placement, and rhetorical structure;
\item source or allusion context supported by an actual window;
\item interpretation, limits, and a project translation.
\end{enumerate}

Apply the layers to the contrast \Latin{superior/inferior}, delayed naming of
the wolf as \Latin{latro}, causal \Latin{siti compulsi}, and the first speech.
The shared stream and spatial arrangement are textual facts; what they imply
about power is interpretation. The phrase \Latin{fauce improba} can support a
claim about characterization, but not an unsupported biography of the poet.

Write a 500-word commentary on four to twelve lines from one received text
reprinted in this packet. Include a scansion or visible-pattern table suited
to that text, a syntax map, project translation, and one paragraph explicitly
stating what the evidence does not establish. If the commentary makes a
source, allusion, or comparative-context claim, use a relevant text reprinted
in the local dossier as its window; otherwise state that no separate source or
allusion claim is being made.

\ifsolutions
\section{Model and points to check}
A concise thesis might be: \English{Phaedrus turns physical elevation into a
dramatic image of unequal power before the wolf invents a verbal grievance.}
Support it from \Latin{superior/inferior}, \Latin{latro}, and the sequence of
arrival, position, and accusation. Do not call the thesis a source reading;
it is a source-grounded inference.
\else
\section{Before continuing}
You should reasonably be able to join witness, meter or visible pattern,
syntax, rhetoric, supported context, translation, interpretation, and limits
in a commentary whose claim classes remain visible. Test the method on a
second excerpt from the packet rather than polishing only the first: identify
its exact local witness and boundary, make the appropriate scansion or visible
pattern record, recover the sentence structure, and support each interpretive
claim with quoted Latin. Keep a useful comparison separate from a dependence
claim and state what the evidence cannot establish. If any layer rests only on
the English rendering, return to the Latin dossier and the P7 commentary
template. Continue when another reader could trace each claim to its proper
kind of printed evidence.
\fi
\end{courseunit}

\begin{courseunit}{P8}
\chapter{Poetry branch: documented verse composition}

Choose one bounded form: four hexameters, two elegiac couplets, one Horatian
line pattern sustained for four lines, or a six- to twelve-line rhythmic hymn
stanza. State the scheme and the intended thought before drafting. Nothing in
the composition may be labeled R.

\section{The composition sequence}

\begin{enumerate}
\item Write the whole thought in plain Latin prose and verify its grammar.
\item Divide the thought into line-sized movements without yet changing words.
\item Mark every natural quantity and every possible positional length or
elision; consult the printed memory sheets.
\item Place secure words in the fixed metrical positions and keep a list of
unplaced necessities.
\item Substitute by meaning and register, not by approximate English gloss.
\item Recheck the written pattern independently from the drafting order and
record every uncertain quantity or division.
\item Recheck syntax across line boundaries and reject any line whose meter
depends on a false form, false quantity, or unintelligible order.
\item Record source influence, exact borrowed phrases, variants considered,
and remaining uncertainty.
\end{enumerate}

\section{Four ledgers}

Maintain a \textbf{quantity ledger} (syllable, mark, reason), a
\textbf{syntax ledger} (word, governor, function), a \textbf{revision ledger}
(old wording, new wording, reason), and a \textbf{source ledger} (feature or
phrase, passage ID, kind of use). A line does not become sound merely because
one ledger looks complete.

After the first finished draft, make three revisions: one for exactness of
thought, one for idiomatic order and diction, and one for pattern. Then write a
prose commentary on your own poem using the P7 layers. The commentary should
identify one attractive feature you removed because it distorted meaning or
source boundaries.

\ifsolutions
\section{Model documentation}
\textbf{P composition note:} \Latin{Versus mei sunt. Formam hexametri ex
Vergilio didici, conversionem distichi ex Ovidio, membrorum aequabilitatem ex
Horatio, petitionum seriem ex sequentia Pentecostes. Nihil tamen auctori
antiquo attribuo quod ipse composui; verba recepta singillatim in tabula
fontium noto.}

\textbf{Project translation.} The verses are my own. I learned the hexameter
form from Virgil, the turn of the couplet from Ovid, balance of members from
Horace, and a sequence of petitions from the Pentecost Sequence. Yet I
attribute to no ancient author what I composed myself; I record received words
individually in the source table.
\else
\section{Before continuing}
You should reasonably be able to plan, pattern, scan, revise, and document a
bounded original poem; reject metrical success bought by grammatical or
semantic failure; and distinguish every source influence from your own Latin.
\fi
\end{courseunit}
